\clearpage\thispagestyle{empty}
\pdfbookmark[0]{Sources for the report on Lindenau’s solar-diameter investigations}{lindenausolar1809-sources}
\begin{center}\textbf{Sources for the report on Lindenau’s solar-diameter investigations}\end{center}
\noindent This AI-translated text follows Carl Friedrich Gauss, \href{https://gdz.sub.uni-goettingen.de/id/PPN236007467}{\emph{Werke}, vol.~VI} (Göttingen, 1874), pp.~517–520. The report appeared unsigned under Göttingen in the \emph{Göttingische gelehrte Anzeigen}, no.~113, 17 July 1809. Its \href{https://images.p.sub.uni-goettingen.de/iiif/image/gdz:PPN319721507_1809:00001125/full/1650,/0/default.jpg}{opening} reports results from Bernhard von Lindenau’s letter of 8 June to Gauss. At the end, the review gives notice of lunar occultations listed by von Wisniewski and passed on by Heinrich Wilhelm Olbers. The collected heading cites pp.~1121–1126. The notice ends on p.~1127 of the journal.
\par\smallskip\noindent On \href{https://images.p.sub.uni-goettingen.de/iiif/image/gdz:PPN319721507_1809:00001129/full/1650,/0/default.jpg}{p.~1125}, the first printing reads $0{,}9631$ and $117/116$ where the collected edition reads $0{,}09631$ and $17/116$. The translations follow the collected text. Gauss’s follow-up of 16 June 1810 (\emph{Werke}, vol.~VI, p.~525) identifies the first of these journal readings as a misprint. On \href{https://images.p.sub.uni-goettingen.de/iiif/image/gdz:PPN319721507_1809:00001131/full/1650,/0/default.jpg}{p.~1127}, the journal prints the occultation table before the final sentence. The collected edition uses two dashes in its place. Translation and typesetting produced by AI.
\clearpage
